
%%%%%%%%%%%%%%%%%%%%%%%%%%%%%%%%%%%%%%
%%            Resultados            %%
%%%%%%%%%%%%%%%%%%%%%%%%%%%%%%%%%%%%%%

\section{Resultados}
\label{sec:resultados}

\subsection{Levantamento de Requisitos}
\label{subsec:levantamento-de-requisitos}

\subsubsection{Requisitos Funcionais}
\label{subsubsec:requisitos-funcionais}

\begin{frame}{Levantamento de Requisitos: Requisitos Funcionais}
	\quad O principal resultado desta etapa foi a realização do levantamento de requisitos funcionais do sistema, que foi realizado através de entrevistas os pesquisadores envolvidos no projeto SmartEnergy \cite{SmartEnergy}. A seguir estão os requisitos funcionais levantados:

	\begin{itemize}
		\item \textbf{Importação de Dados}: A plataforma deverá conseguir importar dados em diferentes formatos, como, por exemplo, arquivos CSV, JSON, XLXS, dentre outros.
		\item \textbf{Pré-processamento de Dados}: A plataforma deverá conseguir realizar o pré-processamento dos dados, como, por exemplo, a normalização dos dados, remoção de dados faltantes, entre outros.
		\item \textbf{Visualização de Dados}: A plataforma deverá conseguir apresentar os dados visualmente em gráficos, em barras, gráficos de dispersão, entre outros.
		\item \textbf{Processamento de Dados}: A plataforma deverá conseguir realizar o processamento dos dados, como, por exemplo, a aplicação de algoritmos de aprendizado de máquina, ex.: regressão linear, regressão logística, entre outros.
		\item \textbf{Exportação de Dados}: A plataforma deverá conseguir exportar os dados em diferentes formatos, como, por exemplo, arquivos CSV, JSON, XLXS, dentre outros.
	\end{itemize}
\end{frame}

\begin{frame}{Requisitos Funcionais II}
	\begin{itemize}
		\item \textbf{Métricas}: A plataforma deverá ser capaz de apresentar métricas de avaliação dos resultados, como, por exemplo, a acurácia, a precisão, a revocação, entre outros.
		\item \textbf{Extensibilidade}: A plataforma deverá ser capaz de ser estendida para novos algoritmos a partir de interfaces de código tipado e aberto.
		\item \textbf{Interface de Usuário}: A plataforma deverá apresentar uma interface de usuário intuitiva e personalizável.
		\item \textbf{Exportação de Modelos}: A plataforma deverá conseguir exportar os modelos de aprendizado de máquina em binários com extensão.pkl e arquivos descritivos com extensão .json.
		\item \textbf{Exportação de Configurações}: A plataforma deverá realizar a exportação as configurações de execução em arquivos com extensão .json.
		\item \textbf{Personalização de Configurações}: A plataforma deverá permitir a personalização das configurações de execução, como, por exemplo, a quantidade de iterações, a quantidade de camadas, entre outros.
		\item \textbf{Extensão da Interface}: A interface da plataforma deverá ser capaz de ser estendida para novas funcionalidades, preservando os estilos e funcionalidades já existentes.
	\end{itemize}
\end{frame}


\subsubsection{Requisitos Não Funcionais}
\label{subsubsec:requisitos-nao-funcionais}


\begin{frame}{Levantamento de Requisitos: Requisitos Não Funcionais I}

	\quad Além dos requisitos funcionais, também foram levantados requisitos não funcionais, que são requisitos que não estão diretamente relacionados com as funcionalidades do sistema, mas que são importantes para a sua qualidade. A seguir estão os requisitos não funcionais levantados:

	\bigskip

	\begin{itemize}
		\item \textbf{Execução Multiplataforma}: A plataforma deverá ser capaz de ser executada em diferentes plataformas, como, por exemplo, Windows, Linux, Mac OS, entre outros.
		\item \textbf{Consumo de Recursos}: A plataforma, principalmente parte que não trata do processamento, deverá conseguir consumir poucos recursos de hardware, como, por exemplo, memória RAM, CPU, entre outros.
		\item \textbf{Confiabilidade}: A plataforma deverá ser capaz de ser executada sem falhas, e em caso de falhas, deverá conseguir ser executada novamente sem a necessidade de intervenção humana.
		\item \textbf{Desempenho}: A plataforma deverá poder ser executada em tempo hábil, sem a atrasar o processo de análise de dados.
	\end{itemize}
\end{frame}

\begin{frame}{Levantamento de Requisitos: Requisitos Não Funcionais II}
	\begin{itemize}
		\item \textbf{Status Operacional}: A plataforma deverá ser capaz de apresentar a situação de operação, como, por exemplo, se está processando os dados, salvando alterações, baixando dados, exportando gráficos ou configurações, entre outros.
		\item \textbf{Log de Execução}: A plataforma deverá conseguir apresentar um log de execução, onde serão apresentados os eventos que ocorreram durante a execução, como, por exemplo, a importação de dados, a exportação de gráficos, a exportação de modelos, entre outros.
		\item \textbf{Isolamento}: A plataforma deverá ser capaz de ser executada em um contêiner, como, por exemplo, Docker, para ser possível a execução em diferentes ambientes.
	\end{itemize}
\end{frame}
